\section{Related Work}
\label{sec:related}

\paragraph{Learning-based next-best-view planning.}
Recent NBV research spans several design philosophies. Projection- and geometry-based planners choose informative views by evaluating explicit hand-crafted utilities over candidate poses, often with strong task priors or object-centric assumptions~\cite{jia2025pbnbv}. Learning-based methods instead estimate view utility directly. GenNBV~\cite{chen2024gennbv} formulates active reconstruction as a generalizable policy-learning problem in free space. VIN-NBV~\cite{frahm2025vinnbv} predicts relative reconstruction improvement for sampled query views. Hestia~\cite{lu2025hestia} introduces a hierarchical policy that predicts look-at points and collision-aware camera placements. These methods show the benefit of learning a view utility, but they still rely on task-specific intermediate states and custom training loops.

\paragraph{Reconstruction-aware and uncertainty-aware active perception.}
Another line of work optimizes NBV through reconstruction quality, consistency, or uncertainty rather than raw coverage. PVP-Recon~\cite{ye2024pvprecon} progressively plans views for sparse-view surface reconstruction using warping consistency. GauSS-MI~\cite{xie2025gaussmi} uses information-theoretic criteria for active 3D reconstruction under Gaussian splatting. Closely related 3DGS-based systems also use rendering quality or uncertainty to drive active mapping and sensing. Compared with these methods, our framework is deliberately simpler in its decision layer: it predicts a single next pose from frozen geometry features and uses one-step differentiable evaluation as supervision.

\paragraph{Geometry foundation models as NBV priors.}
MapAnything~\cite{keetha2025mapanything}, Depth Anything 3~\cite{lin2025depthanything3}, and VGGT~\cite{wang2025vggt} demonstrate that multi-view geometry can be inferred by feed-forward transformers across broad input conditions. Their outputs include point maps, depth, camera-aware features, or geometry tokens that are highly relevant to active perception. However, these models are usually treated as standalone reconstruction systems rather than as frozen perceptual backbones inside an NBV policy. Our work explores the latter design: geometry is delegated to a frozen prior, while the trainable component is restricted to view selection.

\paragraph{Software structure and reproducibility in active vision.}
Active reconstruction papers frequently describe a monolithic method, while the actual implementation mixes rendering, dataset logic, model wrapping, and evaluation in ways that complicate ablations. The repository studied here instead uses a service-oriented architecture with explicit ports, adapters, and workflows. This separation is not merely a software detail. It enables swappable scene encoders, renderer backends, loss functions, and evaluation protocols, which in turn makes the scientific ablations more controlled. Our manuscript treats this decomposition as a first-class part of the method description rather than as incidental engineering.
